\documentclass[conference]{IEEEtran}
\IEEEoverridecommandlockouts
% The preceding line is only needed to identify funding in the first footnote. If that is unneeded, please comment it out.
\usepackage{cite}
\usepackage{amsmath,amssymb,amsfonts}
\usepackage{algorithmic}
\usepackage{graphicx}
\usepackage{textcomp}
\usepackage{xcolor}
\usepackage{float}
\usepackage{array} % put this in your preamble
\usepackage{url}
\def\BibTeX{{\rm B\kern-.05em{\sc i\kern-.025em b}\kern-.08em
    T\kern-.1667em\lower.7ex\hbox{E}\kern-.125emX}}
\begin{document}

\title{Temporal Gait Modifications Under Cognitive Load
Assessed via Monocular Markerless Motion Capture}


\author{
    \IEEEauthorblockN{Hakeem M. Taha}
    \IEEEauthorblockA{\textit{Dept. of Syst. \& Biomed. Eng} \\
    \textit{Cairo University}\\
    hakeem.taha06@eng-st.cu.edu.eg}
    \and
    \IEEEauthorblockN{Seif T. Hegazy}
    \IEEEauthorblockA{\textit{Dept. of Syst. \& Biomed. Eng} \\
    \textit{Cairo University}\\
    seif.hegazy06@eng-st.cu.edu.eg}

    \\[0.5cm]

    \IEEEauthorblockN{Mahmoud M. Mahfouz}
    \IEEEauthorblockA{\textit{Dept. of Syst. \& Biomed. Eng} \\
    \textit{Cairo University}\\
    Mahmoud.Mahfouz07@eng-st.cu.edu.eg}
    \and
    \IEEEauthorblockN{Ziad A. Abdelawad}
    \IEEEauthorblockA{\textit{Dept. of Syst. \& Biomed. Eng} \\
    \textit{Cairo University}\\
    ziad.abd-elawad06@eng-st.cu.edu.eg}
}
\maketitle

\begin{abstract}
Cognitive–motor interference during dual-task walking is
widely recognized as a sensitive indicator of neuromotor
control and early functional decline.
While laboratory-based motion capture systems provide
detailed biomechanical data, the capability of monocular mark-
erless video analysis to detect subtle temporal gait alterations
under cognitive load remains insufficiently characterized.
This study evaluates whether a sagittal-plane monocular
motion capture pipeline can reliably quantify temporal gait
parameters under a cognitive dual-task paradigm in healthy
adults, and compares two gait event detection strategies.
Five participants performed overground walking under single-
task and dual-task (serial subtraction) conditions. A Sports2D
pipeline extracted spatiotemporal parameters including step
time, stance time, swing time, double-support time, and cadence.
A conventional vertical-axis (Y) event detector was compared
against a novel Anterior-Posterior (AP) detector based on
segment-wise linear detrending of horizontal foot trajectories.
Reliability was assessed via Coefficient of Variation (CV) and
Intraclass Correlation Coefficient (ICC). The AP method
achieved CV $<$ 10\% and ICC $>$ 0.6 across all parameters
and conditions, resolving the systematic failures of the vertical
detector. These findings support the viability of monocular
video pipelines for cognitive–motor gait assessment when
appropriate event detection methods are employed.
\end{abstract}
\vspace{3mm}

\begin{IEEEkeywords}
Stochastic Neural Fields, Attractor Dynamics, Physics-Informed Reinforcement Learning (PIRL), Signal-to-Noise Ratio, Gain Modulation, Phase Transitions, Computational Neuroscience, Bifurcation Theory.
\end{IEEEkeywords}

\section{Introduction}
Human gait has traditionally been viewed as a largely automated motor task; however, contemporary biomechanical research recognizes it as a complex cognitive-motor process [1]. When cognitive loads are introduced during walking, such as in dual-task paradigms, significant alterations in spatiotemporal gait parameters often emerge as the brain divides its resources [1]. Historically, quantifying these subtle kinematic shifts required expensive, laboratory-bound 3D motion capture systems or wearable sensor arrays like Inertial Measurement Units (IMUs) [1], [2]. Recently, advancements in computer vision and markerless pose estimation have introduced monocular video analysis as a highly accessible alternative for clinical and remote settings [3], [4]. This literature review synthesizes the foundational concepts of dual-task gait alterations, evaluates the validity of modern 2D video-based motion capture, explores the underlying pose-estimation technologies, and defines the research gap this study aims to address.

\section{Literature Review}
\subsection{Gait Alterations Under Cognitive Load}
Introducing a secondary cognitive task during walking—known as a dual-task paradigm—competes for cortical resources and demonstrably alters gait patterns [1]. To systematically study this phenomenon, Zhou et al. developed the DUO-GAIT dataset, which captures the kinematic effects of cognitive load and fatigue in sixteen healthy young adults [1]. During the study's dual-task condition, participants were required to walk overground for six minutes while continuously subtracting seven from a random four-digit number [1]. The addition of this serial subtraction task significantly impacted spatial and temporal gait parameters, leading to a notable reduction in walking speed ($F(1,15) = 23.75, p < 0.001$) and a decrease in stride length ($F(1,15) = 17.03, p < 0.001$) [1]. Interestingly, Zhou et al. revealed that the gait modifications induced by cognitive dual-tasking were substantially larger than those caused by physical muscle fatigue alone [1]. The highly controlled, multi-sensor IMU data from DUO-GAIT, specifically the Dual-Task Cost (DTC) metrics, establishes a robust, gold-standard reference for the magnitude of temporal gait changes expected under cognitive load [1].

\subsection{Validity of 2D Video-Based Markerless Motion Capture}
While IMU systems and 3D motion capture yield highly accurate kinematic data, they require specialized, expensive equipment that limits their utility in standard clinical environments [1], [2]. Two-dimensional monocular video offers a highly accessible alternative. Stenum et al. rigorously evaluated the validity of 2D video-based pose estimation for human gait analysis by comparing features extracted via OpenPose against simultaneously recorded 3D motion capture in 31 healthy adults [2]. Their findings demonstrated excellent accuracy for spatiotemporal parameters: when calculating individual participant means, the mean absolute errors (MAE) were just 0.01 s for temporal gait parameters (such as step time, stance time, and swing time) and 0.018 m for step length [2]. Because these error margins fall well below published minimal detectable change (MDC) thresholds, the study confirms that 2D monocular video analysis possesses the sensitivity required to detect clinically meaningful gait modifications [2]. However, it is critical to note that their validation was restricted to steady-state, normal walking and did not evaluate the irregular cadences and varied stride intervals typical of dual-task cognitive loading [1], [2].

\subsection{Advancements in Pose Estimation: RTMPose}
The accuracy of any video-based gait analysis is fundamentally dependent on the quality of its underlying pose estimation algorithm [2], [4]. Previous validation studies, such as the one conducted by Stenum et al., relied on OpenPose—an older algorithm that is computationally heavy and less optimized for real-time applications [2], [3]. Recently, the state-of-the-art has advanced significantly with the introduction of RTMPose by Jiang et al. [3]. RTMPose is a high-performance, real-time multi-person pose estimation framework that treats keypoint localization as a classification task [3]. It achieves superior keypoint localization accuracy—reaching 65.3\%. Average Precision on the COCO-WholeBody dataset compared to OpenPose's 44.2\% while maintaining low latency and requiring a fraction of the computational overhead of its predecessors [3].

\subsection{The Sports2D Analytical Framework}
To make advanced pose estimators usable for biomechanical analysis, Pagnon and Kim developed Sports2D, an open-source, end-to-end Python software pipeline [4]. Sports2D integrates the RTMPose engine [3] to automatically detect joint coordinates from standard video or webcam feeds, computing 2D joint and segment angles specifically tailored for sagittal and frontal plane motions [4]. Furthermore, Sports2D applies necessary filtering (such as Butterworth or Gaussian filters) to remove pose estimation jitter and exports the kinematic data into OpenSim-compatible formats, such as .mot and .trc files [4]. It currently stands as the most user-friendly and highly capable open-source tool available for markerless sagittal-plane gait analysis [4].

\subsection{Identified Research Gap and Preliminary Trajectory}
Despite the methodological rigor established by Stenum et al. [2] and the technological capabilities of Sports2D [4], a distinct validation gap remains in the literature. Sports2D has not yet been validated against a reference standard for the extraction of highly variable spatio-temporal gait parameters, nor has any study tested whether a monocular video pipeline can accurately quantify the specific temporal gait alterations induced by cognitive dual-tasking [1], [2], [4]. This study directly addresses that gap by applying the Sports2D framework [4] to the serial-subtraction dual-task paradigm defined by DUO-GAIT [1]. By utilizing statistical reliability metrics—specifically targeting Coefficient of Variation (CV) and Intraclass Correlation Coefficient (ICC) benchmarks—while referencing the stringent accuracy thresholds established by Stenum et al. [2], this research evaluates whether RTMPose-driven video analysis [3] can reliably detect cognitive-motor gait changes. Preliminary video analyses of participants under both single and dual-task loads utilizing this proposed pipeline have already yielded highly promising results, indicating that the framework successfully captures these complex temporal modifications.

\section{Methodology}
\subsection{Data Acquisition Methods}

We capture sagittal-plane video of participants walking at
a self-selected pace using a smartphone camera recording at
120 fps to ensure high temporal resolution (±8.3 ms). The
camera was mounted on a tripod at hip height, positioned
4 meters from a calibrated 4-meter runway to minimize
parallax error [5]. Data are collected under two conditions—
Baseline and Cognitive Load (Serial 7s task)—allowing for
the calculation of the Dual-Task Cost (DTC) via markerless tracking in Sports2D.

\subsection{Software \& Raw Tracked Parameters} 
Sports2D [4] is an open-source Python-based markerless
motion capture pipeline that uses the RTMPose engine [3] to
extract anatomical keypoints from video frames. By detecting
gait events such as heel-strike and toe-off, it measures step
time, stance time, swing time, double-support time, cadence,
and temporal variability.

\subsection{Data Processing Methods}
Video recordings are processed through a custom automated Python pipeline utilizing Sports2D [4] to obtain 2D foot keypoint trajectories across frames. Raw trajectories are first smoothed using a zero-phase 3~Hz Butterworth low-pass filter. Two gait event detection strategies are evaluated and compared.

The \textit{vertical (Y) detector} identifies Heel-Strike (HS) and Toe-Off (TO) events as local minima in the vertical trajectories of heel and toe keypoints, applying physiologically constrained prominence and minimum inter-event distance thresholds (0.70~s between consecutive heel strikes of the same foot).

The \textit{Anterior-Posterior (AP) detector} operates on horizontal (X-axis) foot trajectories, motivated by the principle established by Zeni et al. that AP-axis kinematic extrema reliably correspond to gait events [6]. Within each walking segment (delimited by boundary markers), the algorithm detects walking direction from the mean heel velocity and negates the X-coordinates if the subject walks right-to-left, standardizing forward motion as positive. A degree-1 polynomial is then fit to the segment and subtracted, yielding a residual signal that oscillates around zero. HS events are detected as peaks (heel furthest forward before plant) and TO events as troughs (toe furthest rearward before lift-off) in the detrended signals. Peak prominence is set adaptively as $k \times \sigma$ of the detrended signal ($k = 0.4$), scaling automatically to stride length and pose estimation noise without per-participant tuning. The same 0.70~s minimum inter-event constraint is enforced.

Following event detection, frame indices are converted to timestamps ($t = \text{Frame}_{\text{index}} / \text{fps}$), enabling stride-by-stride computation of stride length, stride time, stance time, swing time, and stance ratio. Outliers are flagged via physiological thresholds (stance ratios outside 0.1–0.95, lengths under 0.2~m) and a z-score pass ($|z| > 3$). Boundary strides are excluded using a defined time margin. Summary statistics (mean and CV) are aggregated across valid strides, and the Dual-Task Cost (DTC) is computed between conditions. The full pipeline is publicly available at \url{https://github.com/Hakeem-Taha-06/DualTaskTemporalGaitAnalysis}.

\subsection{Statistical Analysis}
To systematically establish the reliability and consistency of the pipeline's measurements, we calculate the Coefficient of Variation (CV) for each computed spatiotemporal parameter, as well as for the Dual-Task Cost (DTC) itself. The CV is computed as:
\begin{equation}
CV (\%) = \left( \frac{\sigma}{\mu} \right) \times 100
\end{equation}
where $\sigma$ is the standard deviation and $\mu$ is the mean of the parameter across strides.

The Dual-Task Cost (DTC) represents the relative percentage change in performance between single-task (ST) and dual-task (DT) conditions, calculated as:
\begin{equation}
DTC (\%) = \left( \frac{X_{ST} - X_{DT}}{X_{ST}} \right) \times 100
\end{equation}
where $X_{ST}$ and $X_{DT}$ are the mean parameter values under single-task and dual-task conditions, respectively. Positive DTC values signify a reduction or worsening of the parameter under cognitive load.

Furthermore, the Intraclass Correlation Coefficient (ICC), specifically ICC(1,1), is computed for each of the participants' recorded parameters to robustly assess measurement reliability across trials. To account for unbalanced datasets (varying numbers of strides per subject), the ICC(1,1) is calculated as:
\begin{equation}
ICC(1,1) = \frac{MS_B - MS_W}{MS_B + (n_0 - 1)MS_W}
\end{equation}
where $MS_B$ and $MS_W$ represent the mean squares between and within subjects, respectively, and $n_0$ is the adjusted number of observations per subject. 

The validity of the pipeline is evaluated against predefined statistical thresholds: CV values below 10\% are considered to indicate good measurement consistency (i.e., ``good enough''), and ICC values above 0.6 are considered to represent good reliability.

\newpage

\section{Results}

Results are reported for both the vertical (Y) and AP event detectors. The Y detector failed to meet reliability thresholds for the majority of parameters, while the AP detector achieved acceptable CV and ICC across the full parameter space.

\subsection{Vertical Detector: Spatial vs.\ Temporal Parameter Reliability}
Under the Y detector, no parameter met both reliability thresholds. Aggregate mean CVs for stride length were 14.7\% (left) and 28.3\% (right) in ST, already well above the 10\% threshold. Temporal event-derived parameters were substantially worse: mean stance time CVs reached 40.0\% (left) and 46.1\% (right) in ST, worsening to 54.2\% and 49.7\% in DT. Swing times followed the same pattern, with aggregate mean CVs of 34.2\% (left) and 48.7\% (right) in ST. At the participant level, P03 and P04 drove the worst readings—stance time CVs of 53.5\% and 52.7\% on the left foot—but even the better-performing participants failed to bring aggregate means within threshold. A consistent left–right asymmetry was present, with right-foot CVs equalling or exceeding left-foot values for most temporal parameters, reflecting accumulated tracking instability on both sides rather than a unilateral occlusion effect. ICC(1,1) values were near zero across the full parameter space; the highest observed values were left stride times in ST (0.226) and DT left stride lengths (0.468), neither approaching the 0.6 threshold.

\begin{figure}[H]
\centering
\includegraphics[width=\linewidth]{"cv_comparison.png"}
\caption{Mean CV (\%) per spatiotemporal parameter under the vertical (Y) detector, averaged across left and right feet, for ST and DT conditions. The dashed line indicates the 10\% reliability threshold.}
\label{fig:cv_y}
\end{figure}

\begin{figure}[H]
\centering
\includegraphics[width=\linewidth]{"cv_comparison_ap.png"}
\caption{ICC(1,1) values per spatiotemporal parameter under the vertical (Y) detector, for ST and DT conditions. The dashed line indicates the 0.6 reliability threshold.}
\label{fig:icc_y}
\end{figure}

\subsection{AP Detector: Parameter Reliability}
The AP detector achieved CV $<$ 10\% for every parameter, both feet, and both conditions. Stride times were strongest (ST left: CV = 3.0\%, ICC = 0.864; ST right: CV = 2.7\%, ICC = 0.880). Stance times dropped from aggregate means of 40–46\% (Y detector, ST) to 3.6\% left and 3.4\% right—a roughly 12-fold improvement. Swing times improved from 34–49\% to 3.4–3.9\% in ST. Double-support time was the weakest parameter, peaking at 8.9\% CV in DT, yet remained within threshold. The left–right asymmetry observed under the Y detector was eliminated, consistent with the AP method's independence from vertical keypoint occlusion. ICC(1,1) exceeded 0.6 for all parameters in both conditions except stance ratio in ST (left: 0.566, right: 0.571). Under DT, CVs increased modestly (stride lengths: 3.4–4.3\%; stance times: 4.3–4.9\%), and DT stride length ICC reached 0.929 (left) and 0.908 (right).

\begin{figure}[H]
\centering
\includegraphics[width=\linewidth]{"icc_comparison.png"}
\caption{Mean CV (\%) per spatiotemporal parameter under the AP detector, averaged across left and right feet, for ST and DT conditions. The dashed line indicates the 10\% reliability threshold.}
\label{fig:cv_ap}
\end{figure}

\begin{figure}[H]
\centering
\includegraphics[width=\linewidth]{"icc_comparison_ap.png"}
\caption{ICC(1,1) values per spatiotemporal parameter under the AP detector, for ST and DT conditions. The dashed line indicates the 0.6 reliability threshold.}
\label{fig:icc_ap}
\end{figure}

\section{Conclusion}

This study evaluated two gait event detection strategies within a monocular markerless motion capture pipeline (Sports2D with RTMPose) for quantifying spatiotemporal gait parameters under cognitive dual-task loading in five healthy adults. The vertical (Y) detector, which identifies heel-strike and toe-off events from inflection points in vertical foot trajectories, produced reliable measurements only for stride length. Temporal event-derived parameters failed systematically in 3 of 5 participants, with stance time CVs exceeding 50\% under single-task conditions alone. The failure is mechanistically consistent: vertical inflection point detection requires precise local minima in a signal that collapses into noise under moderate pose estimation jitter, and the problem was compounded by far-side limb occlusion that introduced a systematic left-right asymmetry.

The AP detector resolved these failures entirely. By detrending horizontal foot trajectories segment-wise and applying adaptive prominence thresholds, it produces a clean oscillating signal whose extrema correspond directly to gait events regardless of jitter magnitude or walking direction. All parameters met both the CV $<$ 10\% and ICC $>$ 0.6 thresholds in both conditions, with the single exception of stance ratio ICC in ST (0.57). The left–right asymmetry disappeared, confirming that the Y detector's occlusion sensitivity was the source rather than an intrinsic pipeline limitation. The AP method's robustness across participants who previously produced unusable data demonstrates that the pipeline architecture is sound; the gait event detection stage was the critical bottleneck.

The central limitation that remains is the absence of synchronized ground-truth comparison. Internal consistency metrics confirm that the AP pipeline produces stable, repeatable measurements, but absolute error against a gold-standard IMU reference has not been quantified. DTC values are therefore interpretable as directionally consistent estimates of cognitive-motor effects, but their clinical accuracy cannot yet be established.

\section{Future Work}

Two directions remain most critical. First, validation against a synchronized IMU reference is the primary outstanding requirement. Collecting simultaneous Sports2D video and calibrated IMU data—across both single-task and dual-task conditions—would yield per-stride MAE and minimal detectable change estimates, enabling the pipeline to be characterized for clinical use rather than internal consistency alone. The DUO-GAIT dataset [1] provides an existing reference standard whose paradigm directly matches this study's design and represents a natural validation target.

Second, the AP detector's adaptive prominence constant ($k = 0.4$) was set empirically and has not been optimized across a larger and more diverse population. A systematic evaluation across varying gait speeds, stride lengths, and participant anthropometrics—using ground-truth event labels from synchronized IMU data—would establish whether the default value generalizes or requires per-condition tuning, and would characterize the detector's failure boundaries before deployment in clinical or field settings.

\begin{thebibliography}{00}

\bibitem{Duo}
L. Zhou, E. Fischer, C. M. Brahms, U. Granacher, and B. Arnrich,
"DUO-GAIT: A gait dataset for walking under dual-task and fatigue
conditions with inertial measurement units," Sci. Data, vol. 10, p. 543,
Aug. 2023, doi: 10.1038/s41597-023-02391-w.

\bibitem{2danalysis}
J. Stenum, C. Rossi, and R. T. Roemmich, "Two-dimensional video-
based analysis of human gait using pose estimation," PLOS Comput.
Biol., vol. 17, no. 4, p. e1008935, Apr. 2021, doi: 10.1371/jour-
nal.pcbi.1008935.

\bibitem{Rtm}
T. Jiang, P. Lu, L. Zhang, N. Ma, R. Han, C. Lyu, Y. Li, and
K. Chen, "RTMPose: Real-time multi-person pose estimation based
on MMPose," arXiv:2303.07399, Mar. 2023. [Online]. Available:
https://arxiv.org/abs/2303.07399

\bibitem{sport2d}
D. Pagnon and H. Kim, "Sports2D: Compute 2D human pose and angles
from a video or a webcam," J. Open Source Softw., vol. 9, no. 101, p.
6849, Sep. 2024, doi: 10.21105/joss.06849.

\bibitem{compare}
E. P. Washabaugh, T. A. Shanmugam, R. Ranganathan, and C. Krishnan,
"Comparing the accuracy of open-source pose estimation methods for
measuring gait kinematics," Gait Posture, vol. 97, pp. 188–195, Sep.
2022, doi: 10.1016/j.gaitpost.2022.08.008.

\bibitem{zeni}
J. A. Zeni Jr., J. G. Richards, and J. S. Higginson, "Two simple methods
for determining gait events during treadmill and overground walking using
kinematic data," Gait Posture, vol. 27, no. 4, pp. 710--714, May 2008,
doi: 10.1016/j.gaitpost.2007.07.007.

\bibitem{validity}
K. L. Leung et al., "Validity and reliability of gait speed and knee flexion
estimated by a novel vision-based smartphone application," Sensors, vol.
24, no. 23, p. 7625, Nov. 2024, doi: 10.3390/s24237625.

\end{thebibliography}
\end{document}
